% Generated draft: requires mathematical and source review.
\documentclass[11pt]{article}
\usepackage[margin=25mm]{geometry}
\usepackage{fontspec}
\setmainfont{Latin Modern Roman}
\usepackage{amsmath,amssymb,mathtools,unicode-math}
\setmathfont{Latin Modern Math}
\usepackage{enumitem,needspace,xurl,hyperref}
\setlength{\parindent}{0pt}
\setlength{\parskip}{5pt}
\setlength{\emergencystretch}{3em}
\newcommand{\N}{\mathbb{N}}
\newcommand{\Z}{\mathbb{Z}}
\newcommand{\Q}{\mathbb{Q}}
\newcommand{\R}{\mathbb{R}}
\newcommand{\C}{\mathbb{C}}
\newcommand{\F}{\mathbb{F}}
\newcommand{\A}{\mathbb{A}}
\newcommand{\PP}{\mathbb{P}}
\newcommand{\E}{\mathbb{E}}
\newcommand{\eps}{\varepsilon}
\newcommand{\dd}{\,\mathrm{d}}
\begin{document}
\section*{UnsolvedMath ID: 30001611}
\textbf{Source problem code:} catalogue entry
\subsection{Definitions and mathematical statement}
For orientation, in the usual bounded-operator setting, let \(X\) be a Banach space and let \(T\colon X\to X\) be a bounded linear operator with bounded inverse. The operator \(T\) is called \emph{bicyclic} if there exists \(x\in X\) such that
\[
\overline{\operatorname{span}\{T^k x:k\in\Z\}}=X.
\]
It is called \emph{cyclic} if there exists \(y\in X\) such that
\[
\overline{\operatorname{span}\{T^k y:k=0,1,2,\ldots\}}=X.
\]
Here closures are in the norm topology, spans use the scalar field of \(X\), and operator norms are induced by that norm. These are conventional definitions supplied to explain the terminology; the source record does not specify the ambient space, its scalar field, or the complete hypotheses of the referenced cyclicity theorem.

For positive integers \(n\), the two inverse-growth conditions in the record are
\[
\lim_{n\to\infty}\frac{\log\lVert T^{-n}\rVert}{\sqrt n}=0
\qquad\text{and}\qquad
\log\lVert T^{-n}\rVert=O(\sqrt n).
\]
The second condition means that there exist \(C>0\) and a positive integer \(N\) such that
\[
\bigl|\log\lVert T^{-n}\rVert\bigr|\le C\sqrt n
\qquad\text{for every integer }n\ge N.
\]

\textbf{Conjecture as supplied, with unresolved scope.} The record refers to a theorem asserting that a bicyclic operator satisfying the first growth condition is cyclic, under the remaining hypotheses of that theorem. It conjectures that the first condition cannot be replaced by the second condition while retaining those remaining hypotheses.

More formally, let \(\mathcal H(X,T)\) stand for the full remaining hypotheses of the referenced theorem, which are not supplied. The intended counterexample assertion is
\[
\exists (X,T)\;\Bigl[
\mathcal H(X,T)\ \land\ T\text{ is bicyclic}\ \land\
\log\lVert T^{-n}\rVert=O(\sqrt n)
\ \land\ T\text{ is not cyclic}
\Bigr].
\]
The quantifier ranges over the spaces and operators permitted by the referenced theorem. Because those domains and \(\mathcal H\) are unspecified, this is a conditional formulation of the intended target, not a complete standalone conjecture for all Banach-space operators. The supplied claim of optimality concerns this particular replacement of the growth hypothesis.
\subsection{Short English statement}
The conjecture asks whether the square-root little-\(o\) bound on inverse powers in a bicyclic-to-cyclic theorem can fail to imply cyclicity when weakened to a square-root big-\(O\) bound. A counterexample must satisfy all other hypotheses of that theorem. Those hypotheses and the intended class of spaces are missing from the supplied record, so the precise scope remains unresolved.
\subsection{Sources}
\begin{itemize}
\item Supplied local UnsolvedMath registry record, permanent ID \(30001611\), source problem code \texttt{catalogue entry}, titled ``Optimal Inverse-Growth Conditions for Bicyclic Operator Cyclicity,'' in the Oberwolfach Reports Open Problems collection. The record gives proposed year \(2010\) and does not identify a proposer.
\item Supplied report citation: \emph{Operator Theory and Harmonic Analysis} (2011), Oberwolfach Reports, pp.~2813--2875, DOI \texttt{10.4171/owr/2010/49}. \url{https://doi.org/10.4171/owr/2010/49}. This is also the supplied literature reference labeled ``original problem report.'' The citation year and the year embedded in the DOI are retained as supplied.
\item Supplied literature reference labeled ``Abakumov Atzmon and Grivaux problem statement'': \url{https://ems.press/content/serial-article-files/46308}. This label is retained as source attribution; the registry does not separately identify the conjecture's proposer.
\item Literature assessment checked on \texttt{2026-08-21}: the supplied curation labels the problem open and reports that no proof or counterexample resolving the conjectured replacement was found. This is supplied dated curation derived from the OWR material, not an independently verified literature assessment. The record expressly requests independent verification of time-sensitive status and cited claims.
\item The registry labels the statement \texttt{corrected\_verified} and reports a manual semantic, typographic, and LaTeX review asserting that its cleanup is self-contained and adds no unsupported source-dependent content. These are supplied editorial qualifications, not independent verification of the theorem's hypotheses or mathematical correctness.
\end{itemize}
\end{document}
